\documentclass[a4paper,11pt]{article}
\usepackage[utf8]{inputenc}
\usepackage[T1]{fontenc}
\usepackage[italian]{babel}
\usepackage[top=1.4cm, bottom=1.45cm, left=1.6cm, right=1.6cm, headheight=13pt]{geometry}
\usepackage{tabularx}
\usepackage{array}
\usepackage{enumitem}
\usepackage{xcolor}
\usepackage{fancyhdr}
\usepackage{lastpage}
\usepackage[pdfencoding=auto, hidelinks]{hyperref}

\definecolor{testo}{RGB}{25, 28, 34}
\definecolor{grigio}{RGB}{105, 112, 122}
\definecolor{linea}{RGB}{175, 180, 188}

\color{testo}
\pagestyle{empty}

\newcommand{\campotesto}[3][0]{%
  \raisebox{-1.35mm}{%
    \TextField[
      name=#2,
      width=#3,
      height=5mm,
      align=#1,
      bordercolor={0.68 0.70 0.74},
      backgroundcolor={0.985 0.987 0.992},
      charsize=10pt
    ]{}%
  }%
}

\newcommand{\rispostaaperta}[1]{%
  \TextField[
    name=#1,
    multiline=true,
    donotscroll=true,
    width=\linewidth,
    height=5cm,
    bordercolor={0.68 0.70 0.74},
    backgroundcolor={0.985 0.987 0.992},
    charsize=10pt
  ]{}%
}

\newcommand{\rispostabreve}[1]{%
  \TextField[
    name=#1,
    multiline=true,
    donotscroll=true,
    width=\linewidth,
    height=1cm,
    bordercolor={0.55 0.59 0.65},
    backgroundcolor={0.96 0.97 0.99},
    charsize=10pt
  ]{}%
}

\newcommand{\spaziocorrezioni}{%
  \noindent\fcolorbox{linea}{white}{%
    \begin{minipage}[t][2.3cm][t]{\dimexpr\linewidth-2\fboxsep-2\fboxrule\relax}
      {\scriptsize\textcolor{grigio}{}}
    \end{minipage}%
  }%
}

\newcommand{\grigliapunti}[1]{%
  \vspace{0.25em}
  {\footnotesize\setlength{\tabcolsep}{3.5pt}\renewcommand{\arraystretch}{1.35}%
  \noindent\begin{tabularx}{\linewidth}{|l|*{5}{>{\centering\arraybackslash}X|}>{\centering\arraybackslash}p{1.8cm}|}
    \hline
    \textbf{Valutazione #1} &
    \textbf{Punto 1} \mbox{\textcolor{grigio}{(5\,pt)}} &
    \textbf{Punto 2} \mbox{\textcolor{grigio}{(5\,pt)}} &
    \textbf{Punto 3} \mbox{\textcolor{grigio}{(5\,pt)}} &
    \textbf{Punto 4} \mbox{\textcolor{grigio}{(5\,pt)}} &
    \textbf{Totale} \\
    \hline
    \textbf{Punteggio} & & & & & \hfill $/ 20$ \\
    \hline
  \end{tabularx}}%
}

\begin{document}
\begin{Form}

\begin{center}
    \Huge
    VERIFICA DI INFORMATICA
\end{center}

\vspace{0.2em}
{\small\renewcommand{\arraystretch}{1.65}%
\noindent\begin{tabularx}{\linewidth}{|>{\centering\arraybackslash}X|c|c|}
  \hline
  \textbf{Cognome e Nome:}\;\campotesto{studente_nome}{6.8cm} &
  \textbf{Classe:}\;\campotesto{studente_classe}{1.5cm} &
  \textbf{Data:}\;\campotesto{studente_data}{2.4cm} \\
  \hline
\end{tabularx}}

\vspace{0.2em}
{\renewcommand{\arraystretch}{1.65}%
\noindent
\begin{tabularx}{\linewidth}{|X|c|c|c|}
  \hline
  \textbf{Tabella di valutazione} & \textbf{Definizioni (20 pt)} & \textbf{Domande (80 pt)} & \textbf{Totale / 100} \\
  \hline
  \textcolor{grigio}{Punteggio ottenuto} & & & \\
  \hline
\end{tabularx}}

\bigskip
\noindent{\small\textcolor{testo}{\textbf{ES1. Definisci in modo chiaro, sintetico e tecnicamente corretto i seguenti termini o concetti.} \\
Punti per definizione: \textbf{2 pt} = corretto $\cdot$ \textbf{1 pt} = incompleto/impreciso $\cdot$ \textbf{0 pt} = errato/non svolto.}}

\vspace{0.6em}

\noindent\textit{}

\begin{enumerate}[leftmargin=1.5em, itemsep=0.55em]
  \item \textbf{Definizione} testo definizione: \hfill \textcolor{grigio}{(2\,pt)}\\
  \rispostabreve{def_01}

  \item \textbf{Definizione} testo definizione: \hfill \textcolor{grigio}{(2\,pt)}\\
  \rispostabreve{def_02}

  \item \textbf{Definizione} testo definizione: \hfill \textcolor{grigio}{(2\,pt)}\\
  \rispostabreve{def_03}

  \item \textbf{Definizione} testo definizione: \hfill \textcolor{grigio}{(2\,pt)}\\
  \rispostabreve{def_04}

  \item \textbf{Definizione} testo definizione: \hfill \textcolor{grigio}{(2\,pt)}\\
  \rispostabreve{def_05}

  \item \textbf{Definizione} testo definizione: \hfill \textcolor{grigio}{(2\,pt)}\\
  \rispostabreve{def_06}

  \item \textbf{Definizione} testo definizione: \hfill \textcolor{grigio}{(2\,pt)}\\
  \rispostabreve{def_07}

  \item \textbf{Definizione} testo definizione: \hfill \textcolor{grigio}{(2\,pt)}\\
  \rispostabreve{def_08}

  \item \textbf{Definizione} testo definizione: \hfill \textcolor{grigio}{(2\,pt)}\\
  \rispostabreve{def_09}

  \item \textbf{Definizione} testo definizione: \hfill \textcolor{grigio}{(2\,pt)}\\
  \rispostabreve{def_10}
\end{enumerate}

\vspace{0.35em}
{\footnotesize\setlength{\tabcolsep}{1.5pt}\renewcommand{\arraystretch}{1.35}%
\noindent\begin{tabularx}{\linewidth}{|l|*{10}{>{\centering\arraybackslash}X|}>{\centering\arraybackslash}p{1.4cm}|}
  \hline
  \textbf{Valutazione ES1} &
  \textbf{Def.\,1} &
  \textbf{Def.\,2} &
  \textbf{Def.\,3} &
  \textbf{Def.\,4} &
  \textbf{Def.\,5} &
  \textbf{Def.\,6} &
  \textbf{Def.\,7} &
  \textbf{Def.\,8} &
  \textbf{Def.\,9} &
  \textbf{Def.\,10} &
  \textbf{Totale} \\
  \hline
  \textbf{Punteggio} & & & & & & & & & & & \hfill $/ 20$ \\
  \hline
\end{tabularx}}%
\newpage

\noindent{\small\textcolor{testo}{\textbf{Rispondi in modo chiaro, sintetico e discorsivo, sviluppando tutti i punti guida per ciascuna domanda.} \\
Criteri per punto: \textbf{5 pt} = corretto $\cdot$ \textbf{2,5 pt} = incompleto/impreciso $\cdot$ \textbf{0 pt} = errato/non svolto.}} \\

\noindent\textbf{ES2. Domanda esercizio} \\
Analizza TESTO DELLA DOMANDA sviluppando i seguenti punti:
\begin{enumerate}[leftmargin=1.5em, itemsep=0.08em, topsep=0.15em, label=\arabic*.]
  \item Punto
  \item Punto
  \item Punto
  \item Punto
\end{enumerate}
\vspace{0.3em}
\rispostaaperta{aperta_01}
\spaziocorrezioni
\grigliapunti{ES2}

\vspace{0.6em}

\noindent\textbf{ES3. Domanda esercizio} \\
Analizza TESTO DELLA DOMANDA sviluppando i seguenti punti:
\begin{enumerate}[leftmargin=1.5em, itemsep=0.08em, topsep=0.15em, label=\arabic*.]
  \item Punto
  \item Punto
  \item Punto
  \item Punto
\end{enumerate}
\vspace{0.3em}
\rispostaaperta{aperta_02}
\spaziocorrezioni
\grigliapunti{ES3}

\newpage

\noindent\textbf{ES4. Domanda esercizio} \\
Analizza TESTO DELLA DOMANDA sviluppando i seguenti punti:
\begin{enumerate}[leftmargin=1.5em, itemsep=0.08em, topsep=0.15em, label=\arabic*.]
  \item Punto
  \item Punto
  \item Punto
  \item Punto
\end{enumerate}
\vspace{0.3em}
\rispostaaperta{aperta_03}
\spaziocorrezioni
\grigliapunti{ES4}

\vspace{0.6em}

\noindent\textbf{ES5. Domanda esercizio} \\
Analizza TESTO DELLA DOMANDA sviluppando i seguenti punti:
\begin{enumerate}[leftmargin=1.5em, itemsep=0.08em, topsep=0.15em, label=\arabic*.]
  \item Punto
  \item Punto
  \item Punto
  \item Punto
\end{enumerate}
\vspace{0.3em}
\rispostaaperta{aperta_04}
\spaziocorrezioni
\grigliapunti{ES5}

\end{Form}
\end{document}